\section{Quasi-finite morphisms}
\label{section-quasi-finite}

\noindent
A solid treatment of quasi-finite morphisms is the basis of many developments
further down the road. It will lead to various versions of Zariski's Main
Theorem, behaviour of dimensions of fibres, descent for \'etale morphisms, etc,
etc. Before reading this section it may be a good idea to take a look at
the algebra results in Algebra, Section \ref{algebra-section-quasi-finite}.

\medskip\noindent
Recall that a finite type ring map $R \to A$ is quasi-finite at
a prime $\mathfrak q$ if $\mathfrak q$ defines an isolated point
of its fibre, see Algebra, Definition \ref{algebra-definition-quasi-finite}.

\begin{definition}
\label{definition-quasi-finite}
\begin{reference}
\cite[II Definition 6.2.3]{EGA}
\end{reference}
Let $f : X \to S$ be a morphism of schemes.
\begin{enumerate}
\item We say that $f$ is {\it quasi-finite at a point $x \in X$}
if there exist an affine neighbourhood $\Spec(A) = U \subset X$
of $x$ and an affine open $\Spec(R) = V \subset S$ such that
$f(U) \subset V$, the ring map $R \to A$ is of finite type,
and $R \to A$ is quasi-finite at the prime of $A$ corresponding to $x$
(see above).
\item We say $f$ is {\it locally quasi-finite} if $f$ is
quasi-finite at every point $x$ of $X$.
\item We say that $f$ is {\it quasi-finite} if $f$ is of finite type
and every point $x$ is an isolated point of its fibre.
\end{enumerate}
\end{definition}

\noindent
Trivially, a locally quasi-finite morphism is locally of finite type.
We will see below that a morphism $f$ which is locally of finite type
is quasi-finite at $x$ if and only if $x$ is isolated in its fibre.
Moreover, the set of points at which a morphism is quasi-finite is open;
we will see this in Section \ref{section-Zariski} on Zariski's Main Theorem.

\begin{lemma}
\label{lemma-algebraic-residue-field-extension-closed-point-fibre}
Let $f : X \to S$ be a morphism of schemes.
Let $x \in X$ be a point. Set $s = f(x)$.
If $\kappa(x)/\kappa(s)$
is an algebraic field extension, then
\begin{enumerate}
\item $x$ is a closed point of its fibre, and
\item if in addition $s$ is a closed point of $S$, then
$x$ is a closed point of $X$.
\end{enumerate}
\end{lemma}

\begin{proof}
The second statement follows from the first by elementary topology.
According to Schemes, Lemma \ref{schemes-lemma-fibre-topological}
to prove the first statement
we may replace $X$ by $X_s$ and $S$ by $\Spec(\kappa(s))$.
Thus we may assume that $S = \Spec(k)$ is the spectrum of a field.
In this case, let $\Spec(A) = U \subset X$ be any affine open
containing $x$. The point $x$ corresponds to a prime ideal
$\mathfrak q \subset A$ such that $\kappa(\mathfrak q)/k$
is an algebraic field extension. By
Algebra, Lemma \ref{algebra-lemma-finite-residue-extension-closed}
we see that $\mathfrak q$ is a maximal ideal, i.e., $x \in U$ is a
closed point. Since the affine opens form
a basis of the topology of $X$ we conclude that $\{x\}$ is closed.
\end{proof}

\noindent
The following lemma is a version of the Hilbert Nullstellensatz.

\begin{lemma}
\label{lemma-closed-point-fibre-locally-finite-type}
Let $f : X \to S$ be a morphism of schemes.
Let $x \in X$ be a point. Set $s = f(x)$.
Assume $f$ is locally of finite type.
Then $x$ is a closed point of its fibre
if and only if $\kappa(x)/\kappa(s)$ is
a finite field extension.
\end{lemma}

\begin{proof}
If the extension is finite, then $x$ is a closed point of
the fibre by
Lemma \ref{lemma-algebraic-residue-field-extension-closed-point-fibre}
above. For the converse, assume that $x$ is a closed point
of its fibre. Choose affine opens $\Spec(A) = U \subset X$
and $\Spec(R) = V \subset S$ such that $f(U) \subset V$.
By Lemma \ref{lemma-locally-finite-type-characterize} the ring map
$R \to A$ is of finite type. Let $\mathfrak q \subset A$,
resp.\ $\mathfrak p \subset R$ be the prime ideal corresponding
to $x$, resp.\ $s$. Consider the fibre ring
$\overline{A} = A \otimes_R \kappa(\mathfrak p)$.
Let $\overline{\mathfrak q}$ be the prime of $\overline{A}$
corresponding to $\mathfrak q$. The assumption that $x$
is a closed point of its fibre implies that $\overline{\mathfrak q}$
is a maximal ideal of $\overline{A}$. Since $\overline{A}$
is an algebra of finite type over the field $\kappa(\mathfrak p)$
we see by the Hilbert Nullstellensatz, see
Algebra, Theorem \ref{algebra-theorem-nullstellensatz},
that $\kappa(\overline{\mathfrak q})$ is a finite extension
of $\kappa(\mathfrak p)$.
Since $\kappa(s) = \kappa(\mathfrak p)$ and
$\kappa(x) = \kappa(\mathfrak q) = \kappa(\overline{\mathfrak q})$
we win.
\end{proof}

\begin{lemma}
\label{lemma-base-change-closed-point-fibre-locally-finite-type}
Let $f : X \to S$ be a morphism of schemes which is locally of finite type.
Let $g : S' \to S$ be any morphism. Denote $f' : X' \to S'$ the base change.
If $x' \in X'$ maps to a point $x \in X$ which is closed in $X_{f(x)}$
then $x'$ is closed in $X'_{f'(x')}$.
\end{lemma}

\begin{proof}
The residue field $\kappa(x')$ is a quotient of
$\kappa(f'(x')) \otimes_{\kappa(f(x))} \kappa(x)$, see
Schemes, Lemma \ref{schemes-lemma-points-fibre-product}.
Hence it is a finite extension of $\kappa(f'(x'))$ as
$\kappa(x)$ is a finite extension of $\kappa(f(x))$ by
Lemma \ref{lemma-closed-point-fibre-locally-finite-type}.
Thus we see that $x'$ is closed in its fibre by applying that lemma
one more time.
\end{proof}

\begin{lemma}
\label{lemma-residue-field-quasi-finite}
Let $f : X \to S$ be a morphism of schemes.
Let $x \in X$ be a point. Set $s = f(x)$.
If $f$ is quasi-finite at $x$, then the residue field
extension $\kappa(x)/\kappa(s)$ is finite.
\end{lemma}

\begin{proof}
This is clear from Algebra, Definition \ref{algebra-definition-quasi-finite}.
\end{proof}

\begin{lemma}
\label{lemma-quasi-finite-at-point-characterize}
Let $f : X \to S$ be a morphism of schemes.
Let $x \in X$ be a point. Set $s = f(x)$.
Let $X_s$ be the fibre of $f$ at $s$.
Assume $f$ is locally of finite type.
The following are equivalent:
\begin{enumerate}
\item The morphism $f$ is quasi-finite at $x$.
\item The point $x$ is isolated in $X_s$.
\item The point $x$ is closed in $X_s$
and there is no point $x' \in X_s$, $x' \not = x$
which specializes to $x$.
\item For any pair of affine opens
$\Spec(A) = U \subset X$, $\Spec(R) = V \subset S$ with
$f(U) \subset V$ and $x \in U$ corresponding to $\mathfrak q \subset A$
the ring map $R \to A$ is quasi-finite at $\mathfrak q$.
\end{enumerate}
\end{lemma}

\begin{proof}
Assume $f$ is quasi-finite at $x$. By assumption there exist opens
$U \subset X$, $V \subset S$ such that $f(U) \subset V$, $x \in U$
and $x$ an isolated point of $U_s$. Hence $\{x\} \subset U_s$ is an open
subset. Since $U_s = U \cap X_s \subset X_s$ is also open we conclude
that $\{x\} \subset X_s$ is an open subset also. Thus we conclude
that $x$ is an isolated point of $X_s$.

\medskip\noindent
Note that $X_s$ is a Jacobson scheme by
Lemma \ref{lemma-ubiquity-Jacobson-schemes}
(and
Lemma \ref{lemma-base-change-finite-type}).
If $x$ is isolated in $X_s$, i.e., $\{x\} \subset X_s$ is open,
then $\{x\}$ contains a closed point (by the Jacobson property), hence
$x$ is closed in $X_s$. It is clear that there is no point $x' \in X_s$,
distinct from $x$, specializing to $x$.

\medskip\noindent
Assume that $x$ is closed in $X_s$ and that there is no point $x' \in X_s$,
distinct from $x$, specializing to $x$. Consider a pair of affine opens
$\Spec(A) = U \subset X$, $\Spec(R) = V \subset S$ with
$f(U) \subset V$ and $x \in U$. Let $\mathfrak q \subset A$ correspond to
$x$ and $\mathfrak p \subset R$ correspond to $s$.
By Lemma \ref{lemma-locally-finite-type-characterize} the ring map
$R \to A$ is of finite type. Consider the fibre ring
$\overline{A} = A \otimes_R \kappa(\mathfrak p)$.
Let $\overline{\mathfrak q}$ be the prime of $\overline{A}$ corresponding
to $\mathfrak q$. Since $\Spec(\overline{A})$ is an open subscheme of
the fibre $X_s$ we see that $\overline{q}$ is a maximal ideal
of $\overline{A}$ and that there is no point of $\Spec(\overline{A})$
specializing to $\overline{\mathfrak q}$.
This implies that $\dim(\overline{A}_{\overline{q}}) = 0$.
Hence by
Algebra, Definition \ref{algebra-definition-quasi-finite}
we see that $R \to A$ is quasi-finite at $\mathfrak q$, i.e.,
$X \to S$ is quasi-finite at $x$ by definition.

\medskip\noindent
At this point we have shown conditions (1) -- (3) are all equivalent.
It is clear that (4) implies (1). And it is also clear that
(2) implies (4) since if $x$ is an isolated point of $X_s$
then it is also an isolated point of $U_s$ for any open $U$
which contains it.
\end{proof}

\begin{lemma}
\label{lemma-finite-fibre}
Let $f : X \to S$ be a morphism of schemes.
Let $s \in S$. Assume that
\begin{enumerate}
\item $f$ is locally of finite type, and
\item $f^{-1}(\{s\})$ is a finite set.
\end{enumerate}
Then $X_s$ is a finite discrete topological space, and
$f$ is quasi-finite at each point of $X$ lying over $s$.
\end{lemma}

\begin{proof}
Suppose $T$ is a scheme which (a) is locally of finite type
over a field $k$, and (b) has finitely many points. Then
Lemma \ref{lemma-ubiquity-Jacobson-schemes} shows $T$ is a
Jacobson scheme. A finite Jacobson space is discrete, see
Topology, Lemma \ref{topology-lemma-finite-jacobson}.
Apply this remark to the fibre $X_s$ which is locally of finite type over
$\Spec(\kappa(s))$ to see the first statement. Finally, apply
Lemma \ref{lemma-quasi-finite-at-point-characterize} to see the second.
\end{proof}

\begin{lemma}
\label{lemma-locally-quasi-finite-fibres}
\begin{slogan}
Finite type morphisms with discrete fibers are quasi-finite.
\end{slogan}
Let $f : X \to S$ be a morphism of schemes.
Assume $f$ is locally of finite type.
Then the following are equivalent
\begin{enumerate}
\item $f$ is locally quasi-finite,
\item for every $s \in S$ the fibre $X_s$ is a discrete topological space, and
\item for every morphism $\Spec(k) \to S$ where $k$ is a field
the base change $X_k$ has an underlying discrete topological space.
\end{enumerate}
\end{lemma}

\begin{proof}
It is immediate that (3) implies (2).
Lemma \ref{lemma-quasi-finite-at-point-characterize}
shows that (2) is equivalent to (1).
Assume (2) and let $\Spec(k) \to S$ be as in (3).
Denote $s \in S$ the image  of $\Spec(k) \to S$.
Then $X_k$ is the base change of $X_s$ via
$\Spec(k) \to \Spec(\kappa(s))$. Hence every
point of $X_k$ is closed by
Lemma \ref{lemma-base-change-closed-point-fibre-locally-finite-type}.
As $X_k \to \Spec(k)$ is locally of finite type (by
Lemma \ref{lemma-base-change-finite-type}),
we may apply
Lemma \ref{lemma-quasi-finite-at-point-characterize}
to conclude that every point of $X_k$ is isolated, i.e., $X_k$ has
a discrete underlying topological space.
\end{proof}

\begin{lemma}
\label{lemma-quasi-finite-locally-quasi-compact}
Let $f : X \to S$ be a morphism of schemes.
Then $f$ is quasi-finite if and only if $f$ is
locally quasi-finite and quasi-compact.
\end{lemma}

\begin{proof}
Assume $f$ is quasi-finite. It is quasi-compact by Definition
\ref{definition-finite-type}. Let $x \in X$.
We see that $f$ is quasi-finite at $x$ by
Lemma \ref{lemma-quasi-finite-at-point-characterize}.
Hence $f$ is quasi-compact and locally quasi-finite.

\medskip\noindent
Assume $f$ is quasi-compact and locally quasi-finite.
Then $f$ is of finite type. Let $x \in X$ be a point.
By Lemma \ref{lemma-quasi-finite-at-point-characterize}
we see that $x$ is an isolated point of its fibre.
The lemma is proved.
\end{proof}

\begin{lemma}
\label{lemma-quasi-finite}
Let $f : X \to S$ be a morphism of schemes.
The following are equivalent:
\begin{enumerate}
\item $f$ is quasi-finite, and
\item $f$ is locally of finite type, quasi-compact, and has finite fibres.
\end{enumerate}
\end{lemma}

\begin{proof}
Assume $f$ is quasi-finite. In particular $f$ is locally of finite type
and quasi-compact (since it is of finite type). Let $s \in S$. Since
every $x \in X_s$ is isolated in $X_s$ we see that
$X_s = \bigcup_{x \in X_s} \{x\}$ is an open covering. As $f$
is quasi-compact, the fibre $X_s$ is quasi-compact. Hence we see
that $X_s$ is finite.

\medskip\noindent
Conversely, assume $f$ is locally of finite type, quasi-compact
and has finite fibres. Then it is locally quasi-finite by
Lemma \ref{lemma-finite-fibre}. Hence it is quasi-finite by
Lemma \ref{lemma-quasi-finite-locally-quasi-compact}.
\end{proof}

\noindent
Recall that a ring map $R \to A$ is quasi-finite if it is
of finite type and quasi-finite at {\it all} primes of $A$, see
Algebra, Definition \ref{algebra-definition-quasi-finite}.

\begin{lemma}
\label{lemma-locally-quasi-finite-characterize}
Let $f : X \to S$ be a morphism of schemes.
The following are equivalent
\begin{enumerate}
\item The morphism $f$ is locally quasi-finite.
\item For every pair of affine opens $U \subset X$, $V \subset S$
with $f(U) \subset V$ the ring map
$\mathcal{O}_S(V) \to \mathcal{O}_X(U)$ is quasi-finite.
\item There exists an open covering $S = \bigcup_{j \in J} V_j$
and open coverings $f^{-1}(V_j) = \bigcup_{i \in I_j} U_i$ such
that each of the morphisms $U_i \to V_j$, $j\in J, i\in I_j$
is locally quasi-finite.
\item There exists an affine open covering $S = \bigcup_{j \in J} V_j$
and affine open coverings $f^{-1}(V_j) = \bigcup_{i \in I_j} U_i$ such
that the ring map $\mathcal{O}_S(V_j) \to \mathcal{O}_X(U_i)$ is
quasi-finite, for all $j\in J, i\in I_j$.
\end{enumerate}
Moreover, if $f$ is locally quasi-finite then for
any open subschemes $U \subset X$, $V \subset S$ with $f(U) \subset V$
the restriction $f|_U : U \to V$ is locally quasi-finite.
\end{lemma}

\begin{proof}
For a ring map $R \to A$ let us define
$P(R \to A)$ to mean ``$R \to A$ is quasi-finite''
(see remark above lemma).
We claim that $P$ is a local property of ring maps.
We check conditions (a), (b) and (c) of Definition
\ref{definition-property-local}. In the proof of
Lemma \ref{lemma-locally-finite-type-characterize}
we have seen that (a), (b) and (c) hold for the property
of being ``of finite type''. Note that, for a finite type ring map
$R \to A$, the property $R \to A$ is quasi-finite at $\mathfrak q$
depends only on the local ring $A_{\mathfrak q}$ as an
algebra over $R_{\mathfrak p}$ where $\mathfrak p = R \cap \mathfrak q$
(usual abuse of notation). Using these remarks (a), (b) and (c) of
Definition \ref{definition-property-local} follow immediately.
For example, suppose $R \to A$ is a ring map
such that all of the ring maps $R \to A_{a_i}$ are quasi-finite
for $a_1, \ldots, a_n \in A$ generating the unit ideal.
We conclude that $R \to A$ is of finite type. Also, for any
prime $\mathfrak q \subset A$ the local ring $A_{\mathfrak q}$
is isomorphic as an $R$-algebra to the local ring
$(A_{a_i})_{\mathfrak q_i}$ for some $i$ and some
$\mathfrak q_i \subset A_{a_i}$. Hence we conclude that
$R \to A$ is quasi-finite at $\mathfrak q$.

\medskip\noindent
We conclude that Lemma \ref{lemma-locally-P} applies with $P$
as in the previous paragraph.
Hence it suffices to prove that $f$ is locally quasi-finite is
equivalent to $f$ is locally of type $P$. Since $P(R \to A)$
is ``$R \to A$ is quasi-finite'' which means $R \to A$ is
quasi-finite at every prime of $A$, this follows from
Lemma \ref{lemma-quasi-finite-at-point-characterize}.
\end{proof}

\begin{lemma}
\label{lemma-composition-quasi-finite}
The composition of two morphisms which are locally quasi-finite is
locally quasi-finite. The same is true for quasi-finite morphisms.
\end{lemma}

\begin{proof}
In the proof of Lemma \ref{lemma-locally-quasi-finite-characterize}
we saw that $P = $``quasi-finite'' is a local property of ring maps,
and that a morphism of schemes is locally quasi-finite if and only if
it is locally of type $P$ as in Definition \ref{definition-locally-P}.
Hence the first statement of the lemma follows from
Lemma \ref{lemma-composition-type-P} combined
with the fact that being quasi-finite is a property of ring maps that is
stable under composition, see
Algebra, Lemma \ref{algebra-lemma-quasi-finite-composition}.
By the above, Lemma \ref{lemma-quasi-finite-locally-quasi-compact}
and the fact that compositions of
quasi-compact morphisms are quasi-compact, see
Schemes, Lemma \ref{schemes-lemma-composition-quasi-compact}
we see that the composition of quasi-finite morphisms is
quasi-finite.
\end{proof}

\noindent
We will see later (Lemma \ref{lemma-quasi-finite-points-open})
that the set $U$ of the following lemma is open.

\begin{lemma}
\label{lemma-base-change-quasi-finite}
\begin{slogan}
(Locally) quasi-finite morphisms are stable under base change.
\end{slogan}
Let $f : X \to S$ be a morphism of schemes.
Let $g : S' \to S$ be a morphism of schemes.
Denote $f' : X' \to S'$ the base change of $f$ by $g$
and denote $g' : X' \to X$ the projection.
Assume $X$ is locally of finite type over $S$.
\begin{enumerate}
\item Let $U \subset X$ (resp.\ $U' \subset X'$)
be the set of points where $f$ (resp.\ $f'$) is quasi-finite.
Then $U' = U \times_S S' = (g')^{-1}(U)$.
\item The base change of a locally quasi-finite morphism is
locally quasi-finite.
\item The base change of a quasi-finite morphism is
quasi-finite.
\end{enumerate}
\end{lemma}

\begin{proof}
The first and second assertion follow from the corresponding
algebra result, see
Algebra, Lemma \ref{algebra-lemma-quasi-finite-base-change}
(combined with the fact that $f'$ is also locally of finite type by
Lemma \ref{lemma-base-change-finite-type}).
By the above, Lemma \ref{lemma-quasi-finite-locally-quasi-compact}
and the fact that a base change of a
quasi-compact morphism is quasi-compact, see
Schemes, Lemma \ref{schemes-lemma-quasi-compact-preserved-base-change}
we see that the base change of a quasi-finite morphism
is quasi-finite.
\end{proof}

\begin{lemma}
\label{lemma-quasi-finite-at-a-finite-number-of-points}
Let $f : X \to S$ be a morphism of schemes of finite type.
Let $s \in S$. There are at most finitely many points
of $X$ lying over $s$ at which $f$ is quasi-finite.
\end{lemma}

\begin{proof}
The fibre $X_s$ is a scheme of finite type over a field,
hence Noetherian (Lemma \ref{lemma-finite-type-noetherian}).
Hence the topology on $X_s$ is Noetherian (Properties, Lemma
\ref{properties-lemma-Noetherian-topology})
and can have at most a finite number of isolated points (by elementary
topology). Thus our lemma follows from
Lemma \ref{lemma-quasi-finite-at-point-characterize}.
\end{proof}

\begin{lemma}
\label{lemma-monomorphism-loc-finite-type-loc-quasi-finite}
Let $f : X \to Y$ be a morphism of schemes.
If $f$ is locally of finite type and a monomorphism, then $f$
is separated and locally quasi-finite.
\end{lemma}

\begin{proof}
A monomorphism is separated by Schemes, Lemma
\ref{schemes-lemma-monomorphism-separated}.
A monomorphism is injective, hence we get $f$
is quasi-finite at every $x \in X$ for example by
Lemma \ref{lemma-quasi-finite-at-point-characterize}.
\end{proof}

\begin{lemma}
\label{lemma-immersion-locally-quasi-finite}
Any immersion is locally quasi-finite.
\end{lemma}

\begin{proof}
This is true because an open immersion is a local isomorphism
and a closed immersion is clearly quasi-finite.
\end{proof}

\begin{lemma}
\label{lemma-permanence-quasi-finite}
Let $X \to Y$ be a morphism of schemes over a base scheme $S$.
Let $x \in X$. If $X \to S$ is quasi-finite at $x$, then
$X \to Y$ is quasi-finite at $x$.
If $X$ is locally quasi-finite over $S$, then $X \to Y$
is locally quasi-finite.
\end{lemma}

\begin{proof}
Via Lemma \ref{lemma-locally-quasi-finite-characterize} this translates
into the following algebra
fact: Given ring maps $A \to B \to C$ such that $A \to C$ is
quasi-finite, then $B \to C$ is quasi-finite.
This follows from
Algebra, Lemma \ref{algebra-lemma-four-rings}
with $R = A$, $S = S' = C$ and $R' = B$.
\end{proof}

\begin{lemma}
\label{lemma-quasi-finite-local-source}
Let $f : X \to Y$ and $g : Y \to S$ be morphisms of schemes.
If $f$ is surjective, $g \circ f$ locally quasi-finite, and
$g$ locally of finite type, then $g : Y \to S$  is locally quasi-finite.
\end{lemma}

\begin{proof}
Let $x \in X$ with images $y \in Y$ and $s \in S$.
Since $g \circ f$ is locally quasi-finite by
Lemma \ref{lemma-residue-field-quasi-finite}
the extension $\kappa(x)/\kappa(s)$ is finite.
Hence $\kappa(y)/\kappa(s)$ is finite. Hence
$y$ is a closed point of $Y_s$ by
Lemma \ref{lemma-algebraic-residue-field-extension-closed-point-fibre}.
Since $f$ is surjective, we see that every point of $Y$ is
closed in its fibre over $S$. Thus by
Lemma \ref{lemma-quasi-finite-at-point-characterize}
we conclude that $g$ is quasi-finite at every point.
\end{proof}














